%%%%%%%%%%%%%%%%%%%%%%%%%%%%%%%%%%%%%%%%%%%%%%%%%%%%%%%%%%%%%%%%%%%%%%%%%%%%%%%%
%% MCM/ICM SUMMARY SHEET -- filled in. This file contains the summary page only;
%% no paper body follows.
%%
%% Baseline: .claude/skills/mcm-latex-format/assets/mcm-2027-summary.tex, which is
%% the official COMAP template (corpus/official/MCM-ICM_Summary.tex) with a single
%% correction -- the summary title bar year is 2027 (see corpus/official/INDEX.md
%% section 3.1, conflict 4).
%%
%% Differences from that baseline, and nothing else:
%%   1. \Problem filled in with the chosen problem (was ABCDEF).
%%   2. \Team filled in with the team control number (was 1111111).
%%   3. The official template's red instruction text replaced by the finished
%%      summary. The template itself says "You may delete these instructions"
%%      (MCM-ICM_Summary.tex:54).
%%   4. The baseline's trailing body scaffolding (\clearpage, \pagestyle{fancy},
%%      \newpage, \setcounter{page}{1}, \rhead{Page \thepage\ } and the
%%      "Begin your paper here" placeholder) is omitted, because this file holds
%%      no body: with that scaffolding present, pdfLaTeX emits a second page
%%      carrying only the placeholder and a header. The summary page itself is
%%      unchanged -- \thispagestyle{empty} and \vspace*{-16ex} are kept, so the
%%      summary sheet is still unnumbered (INDEX.md section 2.2.1).
%%   5. The baseline's comment on the \graphicspath line was moved to its own
%%      line. Nothing else on that line changed; a trailing comment after code
%%      is what the mechanical scanner reports as C1 (unescaped percent).
%%
%% Page budget, for the full submission this sheet belongs to (not for this file):
%% the 25-page limit covers the whole submission -- Summary Sheet, Table of
%% Contents, solution, reference list, notes and appendices all count. The only
%% chapter that does not count is "Report on Use of AI", which must come after
%% the 25 pages. See corpus/official/INDEX.md sections 2.1.1 and 2.1.5.
%%
%% This comment block is ASCII-only on purpose.
%%%%%%%%%%%%%%%%%%%%%%%%%%%%%%%%%%%%%%%%%%%%%%%%%%%%%%%%%%%%%%%%%%%%%%%%%%%%%%%%
%%%%%%%%%%%%%%%%%%%%%%%%%%%%%%%%%%%%%%%%
%% MCM/ICM LaTeX Template %%
%% 2027 MCM/ICM           %%
%%%%%%%%%%%%%%%%%%%%%%%%%%%%%%%%%%%%%%%%
\documentclass[12pt]{article}
\usepackage{geometry}
\geometry{left=1in,right=0.75in,top=1in,bottom=1in}

%%%%%%%%%%%%%%%%%%%%%%%%%%%%%%%%%%%%%%%%
% Replace ABCDEF in the next line with your chosen problem
% and replace 1111111 with your Team Control Number
\newcommand{\Problem}{A}
\newcommand{\Team}{2534567}
%%%%%%%%%%%%%%%%%%%%%%%%%%%%%%%%%%%%%%%%

\usepackage{newtxtext}
\usepackage{amsmath,amssymb,amsthm}
\usepackage{newtxmath} % must come after amsXXX

\usepackage[pdftex]{graphicx}
\usepackage{xcolor}
\usepackage{fancyhdr}
\lhead{Team \Team}
\rhead{}
\cfoot{}

\newtheorem{theorem}{Theorem}
\newtheorem{corollary}[theorem]{Corollary}
\newtheorem{lemma}[theorem]{Lemma}
\newtheorem{definition}{Definition}

%%%%%%%%%%%%%%%%%%%%%%%%%%%%%%%%
\begin{document}
% Place your graphic files in the same directory as your main document
\graphicspath{{.}}
\DeclareGraphicsExtensions{.pdf, .jpg, .tif, .png}
\thispagestyle{empty}
\vspace*{-16ex}
\centerline{\begin{tabular}{*3{c}}
	\parbox[t]{0.3\linewidth}{\begin{center}\textbf{Problem Chosen}\\ \Large \textcolor{red}{\Problem}\end{center}}
	& \parbox[t]{0.3\linewidth}{\begin{center}\textbf{2027\\ MCM/ICM\\ Summary Sheet}\end{center}}
	& \parbox[t]{0.3\linewidth}{\begin{center}\textbf{Team Control Number}\\ \Large \textcolor{red}{\Team}\end{center}}	\\
	\hline
\end{tabular}}
%%%%%%%%%%% Begin Summary %%%%%%%%%%%

We build a Stair Wear Model (SWM) that turns a single non-destructive survey of a step
into estimates of how often a stairwell was used, in which direction, and by how many
people at once. The field work needs only a grid of tread-depth readings (straightedge
and dial gauge, or one 3D photogrammetric pass), the step's angle and plan size, and the
stone's hardness; nothing is cored, and two people with hand tools suffice.

The model has two coupled halves. The \emph{Wear Volume Model},
$d(x,y)=T\,N_d\,D(x,y)\,G\,k_m$, sums wear over the tread, so once the wear map $D(x,y)$,
the load $G$ and the coefficient $k_m$ are known, the age $T$ and the daily traffic $N_d$
follow one from the other. The coefficient $k_m$ follows Archard's law $k_m = K d/H$ with
weathering-decayed hardness $H(t)=H_0 e^{-pt}$, so ageing enters the rate; $G$ is a
population-average weight over six demographic groups times $g$, taken here as $G = 700$
N. The \emph{Wear Distribution Model} factorises the map as $D(x,y)=D_X D_Y$ and fits
both marginals with a Gaussian-mixture algorithm (three components across the step, two
along it). Peak counts then read the traffic off the tread: two along-step peaks mean
two-way traffic, and the across-step peak count is how many walked abreast.

Applied to a set of ancient Edinburgh sandstone steps (2.00 m by 0.40 m), the model gives
\begin{itemize}
\setlength{\itemsep}{0pt}
\setlength{\parskip}{0pt}
\setlength{\topsep}{2pt}
\item Frequency of use: $N_d = 261$ person-passes per day; equivalently, at 260 passes
per day the tread was laid $36{,}492$ days ago, about one century.
\item Preferred direction: the along-step marginal splits into two components with an
upward-to-downward ratio of $3:2$, so upward traffic dominated; the fitted impact centres
(0.08 m up, 0.15 m down) both sit near the nosing, the upward one closer, as gait theory
predicts.
\item People abreast: the across-step marginal carries three components at 0.30 m, 1.10 m
and 1.76 m, so three people climbed side by side rather than single file.
\end{itemize}

Four further tests answer the second group of questions. Wear consistency with the record
is scored by a log-ratio index $P_A$ between measured and expected wear matrices; the
paper defines it but reports no value for this stair, so none is claimed here. Age
reliability is a two-sided interval at $\alpha = 0.05$: of the six sandstone steps we
dated, five fell within $5\%$ of the known age, giving $CI \in [0.359, 0.996]$. Repairs
are detected by comparing Brinell hardness across steps, which comes out uniform at 89.45
N/mm$^2$ ($\sigma_H = 1.8$), and by a step-to-step wear scatter $|RS| = 0.03$, below the
$0.05$ similarity threshold; this stair was therefore neither repaired nor rebuilt.
Provenance is tested by comparing the model's time-averaged hardness $H_t = 90.44$
N/mm$^2$ with the literature sandstone value $H_{0At} = 93.18$ N/mm$^2$, an error of
$3.0\%$ -- consistent with sandstone, not with marble or wood.

We also calibrate a traffic-intensity curve mapping the across-step spread $\sigma_x$ onto
the number of people in a typical day,
$N_{\mathrm{pass}} = 0.339 e^{0.13125\sigma_x} + 0.92338$, with $R^2 = 0.9999$;
concentrated footsteps (few people over a long period) and dispersed or multi-modal ones
(many people in a short burst) sit at opposite ends of it. Sensitivity runs show the model
is insensitive to the weathering rate $p$ up to about four centuries but reacts sharply
beyond (raising $p$ from 0.0025 to 0.005 shortens the time to saturation by $64.2\%$), and
it stays robust across four national climates, where a century of simulated wear separates
the extremes by under $0.01\%$. Its limitations are the steady-usage assumption and the
collapse of microclimate into one weathering constant; within those bounds the SWM gives
archaeologists quantitative, non-destructive readings of past traffic.

%%%%%%%%%%% End Summary %%%%%%%%%%%

%%%%%%%%%%%%%%%%%%%%%%%%%%%%%%
\end{document}
\end
